% Alternatives: figure 4 fragment boundary (3), figure 8 spin-weight picture (3), figure 11 covering picture (3).
\documentclass[10pt,border=2mm]{standalone}
\usepackage[T1]{fontenc}
\usepackage{figurestyle}
\input{primitives.tex}
\input{molecules.tex}
\input{numbers.tex}
\begin{document}
\begin{tikzpicture}[plate,boundary/.style={line width=.5pt,draw=PlateInk,dash pattern=on 2.2pt off 1.6pt}]
\path[use as bounding box] (0,0) rectangle (15,26);
% ================================================= figure 4: the fragment boundary
\node[lab,anchor=west] at (.2,25.6) {Figure 4, element: the boundary between the methyl and amino fragments};
\begin{scope}[shift={(2.4,23.3)}]
  \begin{scope}[scale=.8]\draw[boundary,rounded corners=6pt] plot coordinates {\fragloopmethyl} -- cycle; \draw[boundary,rounded corners=6pt] plot coordinates {\fragloopamino} -- cycle;\end{scope}
  \begin{scope}\molsolid\mollabelsfalse\pic[scale=.8] at (0,0) {mol3d-mla-ref};\end{scope}
  \node[sub,anchor=north,text width=4.2cm,align=center] at (0,-1.7) {(i) two dashed contours, each hugging its fragment};
\end{scope}
\begin{scope}[shift={(7.5,23.3)}]
  % (ii) a cut plane between the fragments: the carving idiom; drawn as a thin tinted plate normal to the C-N bond, through its midpoint
  \path[fill=ColNitrogen!14,draw=PlateInk,line width=.5pt] (-.05,-1.75)--(.55,-1.25)--(.55,1.75)--(-.05,1.25)--cycle;
  \begin{scope}\molsolid\mollabelsfalse\pic[scale=.8] at (0,0) {mol3d-mla-ref};\end{scope}
  \node[sub,anchor=north,text width=4.2cm,align=center] at (0,-1.7) {(ii) one cut plane between the fragments: the boundary as a carving cut};
\end{scope}
\begin{scope}[shift={(12.4,23.3)}]
  % (iii) orbit and coset labels after Harter 4.4.1: the methyl hydrogens are one orbit of t, labelled by coset leaders
  \begin{scope}\molsolid\pic[scale=.8] at (0,0) {mol3d-mla-ref};\end{scope}
  \node[sub,anchor=north,text width=4.2cm,align=center] at (0,-1.7) {(iii) no contour: labels alone, fragment by orbit $\{1,2,3,6\}$ and $\{4,5,7\}$, after 4.4.1};
\end{scope}
% ================================================= figure 8: the spin-weight picture
\node[lab,anchor=west] at (.2,20.4) {Figure 8, element: how the 32 spin dimensions pair with the spatial species};
% (i) partitioned bar (current)
\begin{scope}[shift={(.4,18.3)}]
  \pgfmathsetmacro{\u}{.14}
  \path[fill=ColNitrogen!18] (0,0) rectangle ({\u*12},.4); \path[fill=ColOxygen!18] ({\u*12},0) rectangle ({\u*16},.4); \path[fill=ColGold!35] ({\u*16},0) rectangle ({\u*32},.4);
  \foreach \k in {1,...,31} \draw[line width=.25pt,draw=PlateInk!60] ({\u*\k},0)--({\u*\k},.4); \draw[heavy] (0,0) rectangle ({\u*32},.4);
  \draw[heavy] ({\u*1},-.35)--({\u*11},-.35); \draw[heavy] ({\u*12.3},-.35)--({\u*15.7},-.35); \draw[heavy] ({\u*17},-.32)--({\u*31},-.32) ({\u*17},-.38)--({\u*31},-.38);
  \node[sub,anchor=north] at ({\u*6},-.4) {$\mathrm A_1$ 12}; \node[sub,anchor=north] at ({\u*14},-.4) {$\mathrm A_2$ 4}; \node[sub,anchor=north] at ({\u*24},-.4) {$\mathrm E$ 8};
  \node[sub,anchor=north,text width=4.4cm,align=center] at ({\u*16},-1.0) {(i) a 32-unit bar partitioned by spin species, the spatial species as level rules under the block they pair with};
\end{scope}
% (ii) a correlation diagram after 4.2.3: spatial species as level rules on the left, spin species on the right with multiplicities, dashed pairings for even and odd
\begin{scope}[shift={(5.6,16.9)}]
  \node[sub,anchor=south] at (.6,3.1) {$\Gamma_{\mathrm{space}}$}; \node[sub,anchor=south] at (3.4,3.1) {$\Gamma_{\mathrm{spin}}$, copies};
  \foreach \y/\l in {2.6/{\mathrm A_1},1.7/{\mathrm A_2}} {\draw[heavy] (0,\y)--(1.0,\y); \node[sub,anchor=south west] at (0,\y) {$\l$};}
  \draw[heavy] (0,.83)--(1.0,.83) (0,.77)--(1.0,.77); \node[sub,anchor=south west] at (0,.83) {$\mathrm E$};
  \foreach \y/\l/\n in {2.6/{\mathrm A_1}/12,1.7/{\mathrm A_2}/4} {\draw[heavy] (2.9,\y)--(3.9,\y); \node[sub,anchor=south west] at (2.9,\y) {$\l$}; \node[sub,anchor=west] at (4.0,\y) {$\n$};}
  \draw[heavy] (2.9,.83)--(3.9,.83) (2.9,.77)--(3.9,.77); \node[sub,anchor=south west] at (2.9,.83) {$\mathrm E$}; \node[sub,anchor=west] at (4.0,.8) {$8$};
  \draw[fine] (1.0,2.6)--(2.9,2.6) (1.0,1.7)--(2.9,1.7) (1.0,.8)--(2.9,.8);
  \draw[fine,dash pattern=on 2pt off 1.4pt] (1.0,2.6)--(2.9,1.7) (1.0,1.7)--(2.9,2.6);
  \node[sub,anchor=north,text width=4.6cm,align=center] at (2.4,-.1) {(ii) a correlation diagram after 4.2.3: solid lines the even pairing, dashed the odd; the weight is the copy count at right};
\end{scope}
% (iii) dot arrays: the 32 dimensions as dots grouped by species, the spatial species choosing a group
\begin{scope}[shift={(10.9,17.0)}]
  \foreach \i in {0,...,11} {\pgfmathsetmacro{\x}{.3*mod(\i,6)}\pgfmathsetmacro{\y}{2.6-.3*floor(\i/6)}\fill[ColNitrogen] (\x,\y) circle[radius=2.2pt];}
  \node[sub,anchor=west] at (1.9,2.45) {$\mathrm A_1$, $12$};
  \foreach \i in {0,...,3} {\pgfmathsetmacro{\x}{.3*\i}\fill[ColOxygen] (\x,1.7) circle[radius=2.2pt];} \node[sub,anchor=west] at (1.9,1.7) {$\mathrm A_2$, $4$};
  \foreach \i in {0,...,7} {\pgfmathsetmacro{\x}{.45*mod(\i,4)}\pgfmathsetmacro{\y}{1.0-.3*floor(\i/4)}\fill[ColGold] (\x,\y) circle[radius=2.2pt]; \fill[ColGold] (\x+.15,\y) circle[radius=2.2pt];}
  \node[sub,anchor=west] at (1.9,.85) {$\mathrm E$, $8$ pairs};
  \node[sub,anchor=north,text width=4.0cm,align=center] at (1.6,-.1) {(iii) the 32 dimensions as dots grouped by species; counts visible, pairing in words};
\end{scope}
% ================================================= figure 11: the covering picture
\node[lab,anchor=west] at (.2,13.8) {Figure 11, element: twelve sheets over $V$, the two continuations ending apart};
% (i) a stack of oblique plates over a tinted base with lifted paths
\begin{scope}[shift={(.6,7.6)}]
  \path[fill=ColNitrogen!12] (0,0)--(2.6,0)--(3.3,.5)--(.7,.5)--cycle; \draw[heavy] (0,0)--(2.6,0)--(3.3,.5)--(.7,.5)--cycle;
  \node[sub,anchor=west] at (3.4,.25) {$V$};
  \foreach \k in {0,...,5} {\pgfmathsetmacro{\yy}{.9+.72*\k}
    \path[fill=white] (0,\yy)--(2.6,\yy)--(3.3,\yy+.5)--(.7,\yy+.5)--cycle; \draw[fine] (0,\yy)--(2.6,\yy)--(3.3,\yy+.5)--(.7,\yy+.5)--cycle;
    \path[fill=white] (0,\yy+.22)--(2.6,\yy+.22)--(3.3,\yy+.72)--(.7,\yy+.72)--cycle; \draw[fine] (0,\yy+.22)--(2.6,\yy+.22)--(3.3,\yy+.72)--(.7,\yy+.72)--cycle;}
  \coordinate (b) at (1.65,.25); \fill[PlateInk] (b) circle[radius=1.2pt];
  \draw[tpath] (b) .. controls (.5,1.1) and (.4,-.4) .. (b); \draw[upath] (b) .. controls (2.8,-.4) and (2.9,1.1) .. (b);
  \draw[tpath] (1.65,1.15)--(1.65,2.6); \draw[upath] (1.65,2.6)--(1.65,3.05);
  \draw[upath] (1.3,1.15)--(1.3,1.6); \draw[tpath] (1.3,1.6)--(1.3,4.0);
  \node[sub,anchor=north,text width=4.4cm,align=center] at (1.65,-.4) {(i) a stack of plates, two per version, over the tinted base $V$; lifts of the two loops rise to different plates};
\end{scope}
% (ii) the Schreier graph of the 12 sheets: nodes gX, blue t-edges, red dashed b-edges; two paths highlighted
\begin{scope}[shift={(6.0,8.1)}]
  \foreach \k/\l in {0/E,1/t,2/t^2} {\pgfmathsetmacro{\a}{90+120*\k}
    \coordinate (o\k) at ({1.5*cos(\a)},{1.5*sin(\a)}); \coordinate (i\k) at ({.75*cos(\a)},{.75*sin(\a)});}
  \foreach \k in {0,1,2} {\pgfmathsetmacro{\n}{mod(\k+1,3)} \draw[tpath] (o\k) to[bend right=25] (o\n);}
  \foreach \k in {0,1,2} {\pgfmathsetmacro{\n}{mod(\k+2,3)} \draw[tpath] (i\k) to[bend left=25] (i\n);}
  \foreach \k in {0,1,2} \draw[upath,-] (o\k)--(i\k);
  \foreach \k/\l in {0/{X},1/{tX},2/{t^2X}} \node[circle,draw=PlateInk,fill=white,line width=.5pt,inner sep=1pt,font=\fontsize{7}{8}\selectfont] at (o\k) {$\l$};
  \foreach \k/\l in {0/{bX},1/{tbX},2/{t^2bX}} \node[circle,draw=PlateInk,fill=white,line width=.5pt,inner sep=1pt,font=\fontsize{7}{8}\selectfont] at (i\k) {$\l$};
  \node[sub,anchor=north,text width=4.6cm,align=center] at (0,-2.0) {(ii) the sheets as a graph: $t$ in blue, $b$ in red dashed; $t$ then $b$ ends at $tbX$, $b$ then $t$ at $btX=t^2bX$: visibly different nodes (six of twelve shown, the $u$ half alike)};
\end{scope}
% (iii) a helicoid: the loop of V unwinds to a spiral staircase of sheets
\begin{scope}[shift={(11.6,8.0)}]
  \draw[fine] (0,-.6) ellipse[x radius=1.4,y radius=.45];
  \foreach \k in {0,...,11} {\pgfmathsetmacro{\a}{30*\k}\pgfmathsetmacro{\h}{.22*\k}
    \draw[fine] (0,{\h})--({1.4*cos(\a)},{\h+.45*sin(\a)});
    \fill[PlateInk] ({1.4*cos(\a)},{\h+.45*sin(\a)}) circle[radius=1pt];}
  \draw[heavy] plot[domain=0:330,samples=60,variable=\a] ({1.4*cos(\a)},{.22*\a/30+.45*sin(\a)});
  \draw[fine] (0,0)--(0,2.6);
  \node[sub,anchor=north,text width=4.4cm,align=center] at (0,-1.2) {(iii) a helicoid over the loop: one turn downstairs climbs one step; twelve steps before the sheet repeats. Schematic only: the deck group is not cyclic};
\end{scope}
\node[sub,anchor=north west,text width=14cm] at (.3,5.6) {Figure 11 critique: (i) keeps the geometry of a covering (sheets over a base) and the plan's intent; (ii) states the algebra exactly and shows the two endpoints apart without any drawing of sheets; (iii) is wrong in spirit (the covering is not cyclic). Plan: (i) as the main drawing with (ii) as its bookkeeping beneath, replacing the text rows.};
\end{tikzpicture}
\end{document}
